% ============================================================================
%  SECCIÓN 4.4: DATOS SIMULADOS
% ============================================================================

\section{Datos simulados}

En esta etapa del desarrollo, la aplicaci\'on utiliza datos reales almacenados
en Supabase, pero para las funcionalidades que a\'un no est\'an conectadas a la
base de datos se emplean datos temporales dentro de la aplicaci\'on. La
finalidad es comprender c\'omo manipular informaci\'on dentro de la aplicaci\'on
antes de almacenarla en una base de datos.

\subsection{Datos de escenarios (Supabase)}

Los escenarios deportivos se obtienen desde la tabla \textit{scenarios} de
Supabase, que contiene datos reales insertados mediante scripts de
inicializaci\'on (\textit{seed data}). La consulta principal obtiene los
escenarios activos junto con sus im\'agenes, deportes asociados y eventos:

\begin{lstlisting}
const { data, error } = await supabase
  .from('scenarios')
  .select(`
    *,
    scenario_images(url, is_primary, storage_path, display_order),
    scenario_sports(sports(nombre))
  `)
  .eq('estado', 'activo')
  .order('nombre');
\end{lstlisting}

Los datos iniciales incluyen escenarios deportivos de Cochabamba, como
estadios, coliseos y canchas, con informaci\'on de capacidad, direcci\'on y
deportes disponibles.

\subsection{Datos de eventos (temporales)}

Los eventos pr\'oximos de cada escenario se obtienen de la tabla
\textit{events} de Supabase. En esta etapa, los eventos disponibles son
limitados y se utilizan como datos de demostraci\'on para validar la
visualizaci\'on en la pantalla de detalle:

\begin{lstlisting}
interface Event {
  id: string;
  nombre: string;
  descripcion: string | null;
  fecha: string;
  hora: string;
}
\end{lstlisting}

\subsection{Datos de usuario (Supabase Auth)}

La informaci\'on del usuario autenticado se obtiene directamente de
\textit{Supabase Auth}, que gestiona el registro, inicio de sesi\'on y
cierre de sesi\'on. Los metadatos del usuario incluyen el nombre completo y
el rol (\textit{asistente}, \textit{gestor} o \textit{admin}):

\begin{lstlisting}
const { error } = await supabase.auth.signUp({
  email,
  password,
  options: {
    data: {
      full_name: fullName,
      role: 'asistente',
    },
  },
});
\end{lstlisting}

\subsection{Listas din\'amicas}

La pantalla de cat\'alogo implementa listas din\'amicas que responden a los
cambios en los filtros de b\'usqueda. Cuando el usuario escribe un t\'ermino
de b\'usqueda, la lista se actualiza autom\'aticamente para mostrar \\nicamente
los escenarios que coinciden con el criterio:

\begin{center}
\begin{tikzpicture}[
  node distance=0.5cm,
  caja/.style={rectangle, draw=black!60, rounded corners, minimum width=3.5cm,
    minimum height=0.55cm, align=center, font=\small},
  flecha/.style={-Stealth, thick}
]
  \node[caja] (datos) {Datos (Supabase)};
  \node[caja, below=of datos] (estado) {Estado (filtros)};
  \node[caja, below=of estado] (interfaz) {Interfaz (lista)};
  \draw[flecha] (datos) -- (estado);
  \draw[flecha] (estado) -- (interfaz);
\end{tikzpicture}
\end{center}

Esta relaci\'on entre datos, estado e interfaz es fundamental para comprender
c\'omo la aplicaci\'on maneja la informaci\'on de forma reactiva.
